\section{Poczucie szczęścia a miejsce zamieszkania}

% SLAJD 19
\begin{frame}{Wpływ otoczenia na psychikę}
    \begin{itemize}
        \item \textbf{Wybór egzystencjalny:} Decyzja o miejscu zamieszkania to nie tylko logistyka, ale czynnik kształtujący codzienność i relacje.
        \item \textbf{Miasto (Przeciążenie):}
        \begin{itemize}
            \item Permanentne przebodźcowanie (hałas, światła, reklamy).
            \item "Wyścig szczurów" prowadzi do chronicznego stresu i wypalenia.
        \end{itemize}
        \item \textbf{Wieś (Mindfulness):}
        \begin{itemize}
            \item Rytm podyktowany naturą pozwala na regenerację.
            \item \textit{Attention Restoration Theory:} Kontakt z przyrodą odnawia zasoby uwagi i zwiększa kreatywność.
        \end{itemize}
    \end{itemize}
\end{frame}

% SLAJD 20
\begin{frame}{Relacje społeczne i ekonomia}
    \begin{columns}
        \begin{column}{0.5\textwidth}
            \textbf{Paradoks miasta:}
            \begin{itemize}
                \item Tłum ludzi, a jednak samotność.
                \item Relacje transakcyjne i powierzchowne.
                \item Brak poczucia przynależności.
            \end{itemize}
            
            \vspace{0.3cm}
            \textbf{Siła wsi:}
            \begin{itemize}
                \item Model wspólnotowy i kapitał społeczny.
                \item Poczucie bezpieczeństwa dzięki pomocy sąsiedzkiej.
            \end{itemize}
        \end{column}

        \begin{column}{0.5\textwidth}
            \textbf{Aspekt ekonomiczny i zdrowotny:}
            \begin{itemize}
                \item Miasto: Wyższe zarobki, ale ogromne koszty życia (presja finansowa).
                \item Zdrowie: Smog \cite{link3} i hałas vs czyste powietrze.
                \item Wieś: Ten sam budżet = wyższy standard (dom, ogród, pasje).
            \end{itemize}
        \end{column}
    \end{columns}
\end{frame}

% SLAJD 21
\begin{frame}{Zestawienie czynników szczęścia \cite{link2, link3}}
    \begin{table}
        \centering
        \tiny
        \caption{Porównanie wpływu środowiska na dobrostan}
        \begin{tabular}{|l|p{3.5cm}|p{3.5cm}|}
            \hline
            \textbf{Czynnik} & \textbf{Miasto} & \textbf{Wieś / Przedmieścia} \\ \hline
            Poziom stresu & Wysoki, chroniczny (pośpiech, hałas) & Niski, epizodyczny (spokój) \\ \hline
            Więzi społeczne & Anonimowość, powierzchowność & Wspólnota, głębokie relacje \\ \hline
            Środowisko & Smog, hałas & Czyste powietrze, cisza \\ \hline
            Natura & Ograniczona (parki) & Nieograniczona (lasy, pola) \\ \hline
            Finanse & Wysoka presja (drogie życie) & Niższe koszty, wyższy standard \\ \hline
            Bezpieczeństwo & Niższe (anonimowość) & Wyższe (znajomość sąsiadów) \\ \hline
            Tempo życia & Narzucone, szybkie & Zgodne z naturą, wolne \\ \hline
        \end{tabular}
    \end{table}
    
    \vspace{0.2cm}
    \scriptsize
    \textbf{Wniosek:} Obszary wiejskie zapewniają solidniejsze fundamenty dla zdrowia psychicznego i bezpieczeństwa.
\end{frame}

% SLAJD 22
\begin{frame}{Podsumowanie: Świadomy wybór}
    \begin{block}{Definicja szczęścia}
        Szczęście nie zależy od liczby kin czy restauracji, ale od jakości snu, głębokości relacji i poczucia kontroli nad życiem.
    \end{block}

    \begin{itemize}
        \item Migracja na wieś (wspierana pracą zdalną) to świadome odrzucenie ciągłej pogoni za wzrostem.
        \item Wybór równowagi, spokoju i autentyczności.
        \item W kluczowych obszarach dobrostanu, wieś oferuje więcej niż metropolia.
    \end{itemize}
\end{frame}